\section{Sensibilités numériques}
\label{sec-sensibilites}

Les documentations de Guo (§{2.4} de la thèse), de Lisjak et de Tatone et Grasselli consacrent une partie aux paramètres purement numériques : pénalité, taille et orientation du maillage, vitesse de chargement, amortissement. Ce chapitre fait le point pour rockim. Seule la sensibilité à la pénalité a été mesurée en campagne.

\subsection{Pénalité des joints}

Un balayage du facteur de pénalité de 10 à 200 a été joué sur l'onde dans une barre (\cref{sec-V1}). Un modèle de joints en série prédit $M_\text{eff} = M/(1 + \alpha/p_f)$, donc $c_\text{eff}/c = (1 + \alpha/p_f)^{-1/2}$, où $M$ est le module œdométrique. Le modèle suppose $\alpha$ indépendant de $h$ ; comme la longueur de pénalité par défaut est le plus petit diamètre inscrit du maillage (\cref{sec-penalite}), $\alpha$ dépend en fait aussi de la qualité du maillage. La valeur $\alpha = 1{,}21$, calée sur le seul point $p_f = 20$, $h = 4$~mm, a été écrite avant le balayage. La prédiction tient à $0{,}05$ point près de $p_f = 10$ à 200 sur le maillage de 4~mm, écart du même ordre que l'erreur de célérité des éléments finis seuls sur ce maillage ($-0{,}046~\%$, \cref{tab-V1}) ; à $p_f = 20$, l'écart atteint $0{,}34$ point à $h = 2$~mm, alors que l'erreur est indépendante de $h$ à $p_f = 200$. L'ajustement sur les huit points donne $\alpha = 1{,}239$ (\cref{tab-pen,fig-pen}). Le pas de temps décroît comme $1/\sqrt{p_f}$.

\begin{table}[htbp]
\centering
\caption{Retard de célérité de l'onde selon le facteur de pénalité intrinsèque, mesuré et prédit avant calcul.}
\label{tab-pen}
\small
\begin{tabular}{@{}rrrrrr@{}}
\toprule
$p_f$ & $h$ (mm) & célérité mesurée & prédite & erreur $L^2$ & $\Delta t$ (ns) \\
\midrule
10 & 4 & $-5{,}54~\%$ & $-5{,}55~\%$ & $14{,}2~\%$ & $6{,}72$ \\
20 & 4 & $-2{,}92~\%$ & $-2{,}89~\%$ & $7{,}7~\%$ & $4{,}89$ \\
20 & $2{,}83$ & $-3{,}01~\%$ & $-2{,}89~\%$ & $7{,}9~\%$ & $2{,}59$ \\
20 & 2 & $-3{,}23~\%$ & $-2{,}89~\%$ & $8{,}6~\%$ & $1{,}47$ \\
50 & 4 & $-1{,}23~\%$ & $-1{,}19~\%$ & $3{,}3~\%$ & $3{,}14$ \\
100 & 4 & $-0{,}64~\%$ & $-0{,}60~\%$ & $1{,}8~\%$ & $2{,}24$ \\
200 & 4 & $-0{,}35~\%$ & $-0{,}30~\%$ & $1{,}0~\%$ & $1{,}59$ \\
200 & $2{,}83$ & $-0{,}34~\%$ & $-0{,}30~\%$ & $0{,}95~\%$ & $0{,}83$ \\
\bottomrule
\end{tabular}
\end{table}

\begin{figure}[htbp]
\centering
\includegraphics[width=0.6\linewidth]{fig_penalite}
\caption{Retard de célérité de l'onde en fonction du facteur de pénalité intrinsèque, sur deux maillages, avec la prédiction écrite avant le balayage et l'ajustement.}
\label{fig-pen}
\end{figure}

Le choix de la pénalité devient une règle chiffrée : pour un biais de célérité $\varepsilon$, il faut $p_f \geq \alpha/(2\varepsilon) \approx 0{,}62/\varepsilon$, soit 62 pour $1~\%$ et 124 pour $0{,}5~\%$, au prix d'un pas de temps $1{,}8$ ou $2{,}5$ fois plus petit qu'à $p_f = 20$. Les codes de référence fixent la pénalité autrement. Tatone et Grasselli la balaient sur un essai de compression élastique et retiennent la valeur au-delà de laquelle le module ne bouge plus, environ deux ordres de grandeur au-dessus de $E$~\cite{tatone2015}. Yan et al. recommandent au moins 100 fois $E$ d'après une bande sous son poids~\cite{yan2023adaptive}. Guo recommande au contraire $E \leq p_0 \leq 10E$ pour préserver le pas de temps, ce qui correspond, selon la loi mesurée ici et en assimilant les deux définitions de la longueur $h$, à un biais de célérité de 10 à $46~\%$. Le banc de la fissure pressurisée de la campagne applique déjà la règle avec $p_f = 100$.

\subsection{Loi de contact}
\label{sec-sens-contact}

La force par volume de recouvrement de Liu et al. a été comparée au potentiel de Munjiza sur le jumeau grossier de la réplique Kuru (\cref{sec-impact}), dans la configuration de référence (insertion intrinsèque), sur 100~µs, à 4 fils (\cref{tab-potvolume}). Elle réduit le temps de calcul de $36~\%$, mais le pic de la force de contact insert-roche et l'impulsion dépendent du facteur de raideur choisi, et aucun facteur ne reproduit Munjiza à la fois sur le pic et sur l'impulsion. Le facteur $8/3$ n'égale Munjiza qu'en contact face contre face entre tétraèdres égaux ; sur un recouvrement partiel, le rapport des deux forces suit la fraction de face couverte. Un calcul de thèse doit choisir l'une des lois et s'y tenir.

\begin{table}[htbp]
\centering
\caption{Potentiel de Munjiza et force par volume de recouvrement, jumeau grossier Kuru, configuration de référence (insertion intrinsèque), 100~µs, 4 fils.}
\label{tab-potvolume}
\small
\begin{tabular}{@{}lrrrr@{}}
\toprule
loi & temps (s) & pic $F_z$ (N) & impulsion (N\,s) & joints rompus \\
\midrule
Munjiza & 398 & 11\,400 & $0{,}264$ & 57 \\
volume, facteur $8/3$ & 256 & 6\,850 & $0{,}167$ & 56 \\
volume, facteur 4 & 254 & 10\,770 & $0{,}230$ & 36 \\
volume, facteur 5 (défaut) & 255 & 10\,650 & $0{,}223$ & 53 \\
\bottomrule
\end{tabular}
\end{table}

La dispersion du nombre de joints rompus, de 36 à 57, n'a pas été séparée de la variabilité propre à la fragmentation. Ces nombres sont bien inférieurs aux 6\,090 joints rompus du \cref{tab-kuru-insertion} sur le même maillage au même instant : ce tableau a été obtenu le 11 septembre avec le jeu de clés antérieur à la configuration de référence (décharge \cle{origin} sans cliquet). Sur le banc $s = 2{,}5$ à 300~µs, passer au contact de Solidity (raideur $0{,}25\,E$ et relève de naissance \cle{gcBirth = penalty}) réduit la force maximale de $5~\%$ et augmente les ruptures de $13~\%$.

\subsection{Taille du maillage}

Guo a montré que la charge de rupture d'une fissure pressurisée ne converge pas quand la maille passe de 20 à $1{,}25$~mm : elle passe de $4{,}90$ à $3{,}09$~MPa pour $G_f = 50$~J/m$^2$ et de $2{,}63$ à $2{,}10$~MPa pour $G_f = 10$~J/m$^2$. Il en tire la règle d'une maille d'un tiers à un quart de la zone plastique. Lisjak trouve que la résistance en compression converge pour des mailles comprises entre $0{,}4$ et $0{,}15$~mm sur l'argilite d'Opalinus.

rockim n'a pas encore d'étude de convergence en rupture en campagne ; le brésilien adaptatif de la reproduction de Yan ne converge pas quand la maille passe de $1{,}5$ à $0{,}5$~mm (\cref{sec-yan}). Le banc P3.1 reprend le cas de Guo avec des références analytiques (mécanique de la rupture élastique linéaire en largeur finie, $1{,}94$ et $4{,}33$~MPa ; fissure cohésive exacte en plaque infinie, $3{,}33$ et $10{,}15$~MPa) et des critères écrits. Ses essais de mise au point sur les maillages de 20 et 10~mm donnent un amorçage de $1{,}91$ à $1{,}99$~MPa pour $G_f = 10$~J/m$^2$ en insertion intrinsèque, proche de la mécanique de la rupture et sous la solution cohésive en plaque infinie ; ces deux références ne font qu'encadrer, et la référence pertinente pour un modèle cohésif, une solution cohésive en largeur finie, manque ; l'insertion adaptative n'insère aucun joint à 20~mm et amorce à $4{,}48$~MPa à 10~mm. La variante la plus fine est estimée entre 71 et 114~h de calcul. Sur l'impact, la réplique St Anne n'a tourné que sur un maillage, et le jumeau Kuru à maillage grossier n'est pas une étude de résolution, son piston étant plus léger.

Le guide d'utilisation du code (\cle{GUIDE_rockim.md}) pose en amont une condition de faisabilité, $2h < \ell_{cz} = E G_I/f_t^2 < a$, où $a = \sqrt{R\delta}$ est le rayon de contact de Hertz : la borne basse demande deux éléments dans la zone cohésive, la borne haute une zone cohésive contenue dans la zone chargée. Elle s'appuie sur un contre-exemple, un calcul de percussion où la zone cohésive (35~mm) dépassait largement le rayon de contact ($1{,}45$~mm) et n'a rompu que 4 joints sur 158\,423, et sur un seul exemple favorable ; elle n'est pas démontrée. Les deux répliques d'impact du \cref{sec-impact} la violent : avec leurs paramètres, $\ell_{cz}$ vaut 14~mm pour St Anne et 25~mm pour Kuru, plusieurs fois le rayon de contact. Soit la borne haute est trop stricte, soit ces répliques sont hors de son domaine ; le rapport ne tranche pas, et l'absence d'étoile radiale sur Kuru pourrait en relever.

\subsection{Schéma d'insertion}
\label{sec-insertion}

Le choix entre insertion intrinsèque et adaptative n'est pas un réglage numérique neutre : il change la réponse en rupture. Sur le jumeau grossier Kuru, à 100~µs, l'insertion intrinsèque rompt 6\,090 joints et forme une zone broyée de $23{,}6$~mm de profondeur sous l'insert ; l'insertion adaptative en rompt 27 et ne forme qu'une peau de $2{,}3$~mm, avec des tétraèdres restés élastiques jusqu'à 1\,359~MPa (\cref{tab-kuru-insertion}). Évaluer le critère d'insertion sur le maximum des points de la face plutôt que sur leur moyenne (\cle{facetAverage = max}) multiplie les ruptures par 15 et porte la profondeur à $8{,}3$~mm ; un écart subsiste avec l'intrinsèque, 6\,090 contre 412 joints rompus, attribué à l'endommagement diffus accumulé dans les joints de pénalité, propriété de la discrétisation sur laquelle Yang et al. ont calé leurs paramètres. L'insertion adaptative double environ le pas de temps ($6{,}32$ contre $2{,}83$~ns). Sur un tunnel 2D, elle fragmente à l'échelle de l'élément ($75{,}7~\%$ de blocs formés d'un seul élément) et propage moins les fissures. Yan et al. comparent les deux schémas à paramètres égaux sur l'essai brésilien et la compression~\cite{yan2023adaptive} ; rockim reproduit cette comparaison au \cref{sec-yan}, où le plateau intrinsèque reste 13 à $16~\%$ sous l'insertion adaptative. La configuration de référence de l'impact reste intrinsèque, comme Solidity, et le choix fait partie des questions de la thèse.

\subsection{Vitesse de chargement et amortissement}

Lisjak et Tatone montrent que la résistance devient indépendante de la vitesse des plateaux sous environ $0{,}25$~m/s et de l'amortissement au-delà de $0{,}01$ fois l'amortissement critique. Guo trouve un pic de flexion stable sous $0{,}01$~m/s. Ces deux sensibilités n'ont pas de banc de campagne ; elles n'ont été mesurées qu'en marge du tunnel (\cref{sec-tunnel-duree}) et de la calibration (\cref{sec-calib}). La vitesse fait partie des critères du banc de la fissure en aile, préparé ; l'amortissement n'a pas de banc. Un constat ancien en marque l'enjeu : sur un cas mal dimensionné, l'amortissement numérique a absorbé $12~\%$ de l'énergie contre $0{,}55~\%$ pour la fissuration, et sur un tunnel 2D, passer l'amortissement local de $0{,}15$ à $0{,}7$ divise l'endommagement par 4.
